\documentclass{subfiles}
\begin{document}
    To give the most accessible idea of what each of groups, rings and fields are,
        we can think of this as sets on which some operations are defined; 
        in addition, such operations must satisfy some properties.

    Formally speaking, given a set \(G\) and a binary operation \(* : G \to G\)\footnotemark,
        we say that \((G, *)\) forms a group if all of the following hold.
        \begin{enumerate}
            \item There exists an element \(\unit_{G}\) such that \( a * \unit_{G} = \unit_{G} * a = a,
                \forall a \in G\). In other words, \(G\) must have a neutral element.

            \item For all \(a \in G\) there must exists \(a^{-1}\) such that 
                \(a * a^{-1} = a^{-1} * a = \unit_{G}\). 
                Meaning that each element must have an inverse.

            \item The operation \(*\) must be associative.
                That is, for all \(a, b, c \in G\) it holds \( a * (b * c) = (a * b) * c\).
        \end{enumerate}
        Additionally, if \(*\) is commutative, i.e., \(a * b = b * a\),
        we say that \((G, *)\) is an abelian (or commutative) group.

    \footnotetext{
        Note that the use of \(*\) for the binary operation should not fool the reader in 
        to thinking it as the well known multiplication. 
        Instead, it is a placeholder for any well defined operation.
    }

    As for sets, groups may contain some smaller groups, to which we refer to as sub-groups.
        Formally, given \((G, *)\) a group, we say that \(H \subseteq G\) forms a group under 
        \(*\) if \(\unit_G \in H\) and \((H, *)\) forms a group.
        
    Beyond the set-theoretic relation, groups and subgroups are connected by a much more 
        foundamental relation, that regarding their size. This result,
        known in the literature as \emph{Lagrange's theorem for groups},
        states that a finite group \(G\), i.e. a group with finitely many elements,
        can only have subgroups \(H\) such that the number of elements in \(H\) 
        divides the number of elements in \(G\).
        We Formally state Lagrange's theorem below.
        \begin{theorem*}[Lagrange's theorem for groups]
            Let \((G, *)\) be a group, with \(\abs{G} < \infty\). Then, 
                any subgroups \((H, *)\) of \((G, *)\) is such that
                \[
                    \abs{H} \vert \abs{G}.
                \]
        \end{theorem*}

    Although groups have some nice properties as the whole group-theory shows,
        more often than not we need something more powerful.
        To provide for this need, some more advanced structures were developed,
        the simplest of which is known as \emph{ring}.

    To undestand rings, we can see them as an extention of groups by an additional operation.
        More precisely, we say that \(R, +, *\) is a ring if 
        \begin{enumerate}
            \item \((R, +)\) form an abelian group.
            \item \(*\) is associative. That is, for all \(a, b, c \in R\) it holds 
                \(a * (b * c) = (a * b) * c\).
            \item \(*\) distributes over \(+\). In other words,
                if for all \(a, b, c \in R\) we have that \(a * (b + c) = a * b + a * c\).
        \end{enumerate}
        
    Rings can sometimes satisfy some additional properties,
        depending on which they take different names.
        These additional properties are the following.
        \begin{enumerate}
            \item There exists a neutral element for \(*\).
            \item \(*\) is commutative.
            \item Each element in \(R\) as inverse with respect to \(*\).
        \end{enumerate}
        A ring with all these additional properties is called \emph{field}.
        In these notes we denote a generic field by either \(\F\) or \(\K\),
        while finite fields will be denoted by \(\gf(p)\), 
        where \(p\) is usually a prime.
\end{document}
